\documentclass[a4paper, 11pt]{report}



\usepackage[utf8]{inputenc} 
\usepackage[T1]{fontenc}
\usepackage{lmodern}
\usepackage{graphicx}
\usepackage[french]{babel}
\usepackage{color}
%\usepackage{fullpage}
\usepackage{array}
\usepackage[tight]{shorttoc}
\usepackage[toc,page]{appendix} 
\usepackage{makeidx} 
\usepackage{titlesec} % pour enlever "chapitre"


\definecolor{gris}{gray}{0.45}

\newcommand{\rouge}[1]{\textcolor{red}{#1}}
\newcommand{\rem}[1]{\textcolor{blue}{\emph{Remarque: \\} #1}}
\newcommand{\exl}[1]{\textcolor{gris}{\emph{Exemple: } #1}}
\newcommand{\exc}[1]{\textcolor{gris}{(\emph{Ex:} #1)}}
\newcommand{\cod}[1]{\textcolor{gris}{\emph{#1}}} 
\newcommand{\att}[1]{\textcolor{red}{\emph{\\ Attention: \\} #1 \\}}

\usepackage{listings}
\definecolor{dkgreen}{rgb}{0,0.6,0}
\definecolor{gray}{rgb}{0.5,0.5,0.5}
\definecolor{mauve}{rgb}{0.58,0,0.82}
\definecolor{red}{rgb}{1,0,0}

\newcommand{\lstconfig}[1]{
	\lstset{
	  language=#1,				      % the language of the code
	  basicstyle=\footnotesize,	      % the size of the fonts that are used for the code
	  numbers=left,				      % where to put the line-numbers
	  numberstyle=\footnotesize,	  % the size of the fonts that are used for the line-numbers
	  stepnumber=1,				      % the step between two line-numbers. If it's 1, each line 
									  % will be numbered
	  numbersep=5pt,				  % how far the line-numbers are from the code
	  backgroundcolor=\color{white},  % choose the background color. You must add \usepackage{color}
	  showspaces=false,			      % show spaces adding particular underscores
	  showstringspaces=false,		  % underline spaces within strings
	  showtabs=false,				  % show tabs within strings adding particular underscores
	  frame=single,				      % adds a frame around the code
	  tabsize=2,					  % sets default tabsize to 2 spaces
	  captionpos=b,				      % sets the caption-position to bottom
	  breaklines=true,				  % sets automatic line breaking
	  breakatwhitespace=false,		  % sets if automatic breaks should only happen at whitespace
	  title=\lstname,	   			  % show the filename of files included with \lstinputlisting;
									  % also try caption instead of title
	  numberstyle=\tiny\color{gray},  % line number style
	  keywordstyle=\color{blue},	  % keyword style
	  commentstyle=\color{dkgreen}\textit,   % comment style
	  stringstyle=\color{mauve}\textbf,	  % string literal style
	}
}

	%\title{\vspace{5cm}Nascar \\ Rapport de projet \\ LO41 \\ \ \\}
	%\date{automne 2012\\ \ \\ \vspace*{3cm}
	%\includegraphics[width=4cm]{logo_utbm.png}
	%}
	%\author{Paul Locatelli - Pierre Rognon \\ \ \\ \ \\ Université de Technologies de Belfort-Montbéliard\\ \ \\}
	
	\titleformat{\chapter}[hang]{\bf\huge}{\thechapter}{2pc}{}     

\begin{document}

\makeatletter
	\def\clap#1{\hbox to 0pt{\hss #1\hss}}%
	\def\ligne#1{%
	\hbox to \hsize{%
	\vbox{\centering #1}}}%
	\def\haut#1#2#3{%
	\hbox to \hsize{%
	\rlap{\vtop{\raggedright #1}}%
	\hss
	\clap{\vtop{\centering #2}}%
	\hss
	\llap{\vtop{\raggedleft #3}}}}%
	\def\bas#1#2#3{%
	\hbox to \hsize{%
	\rlap{\vbox{\raggedright #1}}%
	\hss
	\clap{\vbox{\centering #2}}%
	\hss
	\llap{\vbox{\raggedleft #3}}}}%
	\def\maketitle{%
	\thispagestyle{empty}\vbox to \vsize{%
	\haut{}{\@blurb}{}
	\vfill
	\vspace{1cm}
\begin{flushleft}
	%\usefont{OT1}{ptm}{m}{n}
	\huge \@title
\end{flushleft}
	\par
	\hrule height 4pt
	\par
\begin{flushright}
	%\usefont{OT1}{phv}{m}{n}
	\Large \@author
	\par
\end{flushright}
	\vspace{1cm}
	\vfill
	\vfill

\begin{center}
	\includegraphics[width=5cm]{logo_UTBM.jpg}
\end{center}

\bas{}{Automne 2012}{}
}%
\cleardoublepage
}
\def\date#1{\def\@date{#1}}
\def\author#1{\def\@author{#1}}
\def\title#1{\def\@title{#1}}
\def\location#1{\def\@location{#1}}
\def\blurb#1{\def\@blurb{#1}}
\date{\today}
\author{}
\title{}

% informations
\location{Belfort}\blurb{}
\makeatother
\title{Projet de résolution d'un tangram}
\author{\small{Adrien \bsc{Berthet} Paul \bsc{Locatelli} et Pierre \bsc{Rognon}}}
\blurb{%
	\textbf{IA41 - Concepts fondamentaux en Intelligence Artificielle et langages dédiés}\\
	Université de Technologie de Belfort-Montbéliard
}% 


	
	\maketitle
	
	\newpage
	
	\shorttoc{Sommaire}{0}
	
	
	\chapter*{Introduction}
	\addcontentsline{toc}{chapter}{Introduction}
	
		Pour ce projet, le sujet abordant le casse-tête du tangram (numéro 6) a été retenu. Le choix de ce sujet s'est fait sur l'intérêt de ce jeu. En effet, s'il apparaît assez aisé de résoudre des modèles tels que le classique "carré" que chacun connaît, d'autres modèles sont beaucoup plus complexes. De plus, c'est un test utilisé dans de nombreux tests afin de déceler d'éventuelles déficiences chez les enfants tels que le WISC (Wechsler Intelligence Scale for Children). C'est donc un problème d'Intelligence Artificielle fondamental qui sera abordé dans ce projet et dont la résolution sera tentée.\\
L'analyse de ce problème constitue une partie prépondérante dans la tentative de résolution du tangram. Environ 6000 configurations différentes sont connues aujourd'hui. De nombreux modèles ont donc dû être envisagés afin de couvrir l'ensemble des cas possibles. Une seconde phase qui s'avère aussi complexe est le choix des différentes représentations informatiques du problème. De nombreux traitements ont été effectués pour mettre à bien cette partie. Quelques exemples qui ont été utilisés comme support durant toute la durée du projet seront présentés afin d'indiquer leur solution. Enfin, la résolution par le logiciel crée sera abordée pour ces exemples. Les limites de ce problème et l'état d'avancement du projet seront indiqués enfin.


	\newpage

	\chapter{Analyse du casse-t\^ete}
	
	Une première analyse du tangram a permis de mettre à jour plusieurs sous-problèmes. Deux problèmes d'ordre mathématique ainsi qu'un problème d'intelligence artificielle pure ont ainsi été isolés. 
	
		\section{Soustraction d'une forme au modèle}
		
		Le second problème mathématique que la résolution exigeait de résoudre est la soustraction d'une forme au modèle. En effet, la solution la plus simple résidait dans le fait de placer une forme dans le modèle pour ensuite renvoyer un nouveau problème et recommencer le casse-t\^ete la forme placée en moins avec le nouveau modèle.
	
	
	\newpage
	
	\chapter{Représentation informatique du problème}
	
	Le passage de la phase d'analyse à la phase "informatique" du pour la résolution du tangram s'avère beaucoup plus complexe que ce qui avait été prévu. Cependant, des représentations précises ont été adoptées.
	
	\chapter{Résultats obtenus}
	
	De nombreux problèmes ayant été rencontrés durant le projet, les résultats obtenus n'indiquent pas une résolution totale de différents modèles. Ceci est d\^u au fait des différentes difficultés apparues sur les algorithmes mathématiques touchant à la géométrie.
	
	
		
	\newpage	
		
	\tableofcontents
		
\end{document}